\documentclass[11pt,a4paper]{article}

% --- Packages ---
\usepackage[utf8]{inputenc}
\usepackage[T1]{fontenc}
\usepackage[french]{babel}
\usepackage{geometry}
\geometry{
    a4paper,
    total={170mm,257mm},
    left=20mm,
    top=25mm,
    bottom=25mm,
    right=20mm
}
\usepackage{graphicx}
\usepackage{booktabs}
\usepackage{xcolor}
\usepackage{hyperref}
\usepackage{listings}
\usepackage{enumitem}
\usepackage{fancyhdr}
\usepackage{titlesec}
\usepackage{tabularx}
\usepackage{tcolorbox}
\usepackage{amsmath}

% --- Custom Colors (Matching WEJHA Design System) ---
\definecolor{wejhaprimary}{HTML}{0F6E56}       % Primary Green
\definecolor{wejhasecondary}{HTML}{1D9E75}     % Secondary Green
\definecolor{wejhaaccent}{HTML}{BA7517}        % Accent Amber
\definecolor{wejhadark}{HTML}{05281E}          % Deep Dark Green/Black
\definecolor{wejhabg}{HTML}{F8F5EF}            % Light Sand Background
\definecolor{lightgrey}{HTML}{F3F4F6}          % Light Grey for tables

% --- Hyperref Setup ---
\hypersetup{
    colorlinks=true,
    linkcolor=wejhaprimary,
    filecolor=wejhaaccent,      
    urlcolor=wejhasecondary,
    pdftitle={Cahier des Charges - WEJHA},
    pdfpagemode=FullScreen,
}

% --- Header and Footer Setup ---
\pagestyle{fancy}
\fancyhf{}
\fancyhead[L]{\textcolor{wejhaprimary}{\textbf{WEJHA}} -- Plateforme Touristique Algérienne}
\fancyhead[R]{\textcolor{gray}{\leftmark}}
\fancyfoot[C]{\thepage}
\fancyfoot[R]{\textcolor{wejhaprimary}{MVP Specifications}}
\renewcommand{\headrulewidth}{0.4pt}
\renewcommand{\footrulewidth}{0.4pt}

% --- Title Section Customization ---
\titleformat{\section}
  {\color{wejhaprimary}\normalfont\Large\bfseries}
  {\thesection}{1em}{}
\titleformat{\subsection}
  {\color{wejhasecondary}\normalfont\large\bfseries}
  {\thesubsection}{1em}{}
\titleformat{\subsubsection}
  {\color{wejhaaccent}\normalfont\normalsize\bfseries}
  {\thesubsubsection}{1em}{}

% --- Code Listing Customization ---
\lstdefinelanguage{TypeScript}{
  keywords={abstract, any, as, boolean, break, case, catch, class, console, const, continue, debugger, declare, default, delete, do, else, enum, export, extends, false, finally, for, from, function, get, if, implements, import, in, instanceof, interface, let, module, new, null, number, of, package, private, protected, public, readonly, require, return, set, static, string, super, switch, this, throw, true, try, type, typeof, var, void, while, with, yield},
  morecomment=[l]{//},
  morecomment=[s]{/*}{*/},
  morestring=[b]',
  morestring=[b]"
}

\lstset{
    language=TypeScript,
    backgroundcolor=\color{lightgrey},
    basicstyle=\ttfamily\small,
    breakatwhitespace=false,         
    breaklines=true,                 
    captionpos=b,                    
    keepspaces=true,                 
    numbers=left,                    
    numbersep=5pt,                  
    showspaces=false,                
    showstringspaces=false,
    showtabs=false,                  
    tabsize=2,
    commentstyle=\color{gray},
    keywordstyle=\color{wejhaprimary}\bfseries,
    stringstyle=\color{wejhaaccent}
}

% --- Document ---
\begin{document}

% --- Page de Garde ---
\begin{titlepage}
    \centering
    \vspace*{1.5cm}
    
    % Logo Placeholder or Title
    \begin{tcolorbox}[colback=wejhaprimary,colframe=wejhaprimary,arc=15pt,halign=center,width=4cm]
        \textcolor{white}{\Huge\textbf{W}}
    \end{tcolorbox}
    
    \vspace{1cm}
    \textcolor{wejhaprimary}{\Huge \textbf{WEJHA}}\\
    \vspace{0.4cm}
    \textcolor{gray}{\Large Spécifications Fonctionnelles et Techniques}\\
    \vspace{0.2cm}
    \textcolor{wejhaaccent}{\large Version MVP -- Algérie}
    
    \vspace{2.5cm}
    
    \begin{tcolorbox}[colback=wejhabg, colframe=wejhaprimary, title=Informations du Projet, arc=8pt]
        \begin{tabular}{ll}
            \textbf{Auteur :} & Équipe de Développement WEJHA \\
            \textbf{Date :} & \today \\
            \textbf{Statut :} & Prêt pour Copy-Paste \\
            \textbf{Technologies clés :} & React, TypeScript, Vite, Google Gemini API \\
            \textbf{Version du document :} & 1.1.0
        \end{tabular}
    \end{tcolorbox}
    
    \vspace{3cm}
    
    \textcolor{wejhadark}{\large\textit{«~Votre direction ultime pour explorer la richesse du patrimoine et des sites naturels en Algérie.~»}}
    
    \vfill
    
    \textcolor{gray}{\small Document confidentiel -- Destiné à l'équipe technique}
\end{titlepage}

\newpage
\tableofcontents
\newpage

% --- Section 1: Introduction ---
\section{Présentation du Projet}
\subsection{Contexte et Objectifs}
L'Algérie dispose d'un patrimoine culturel, historique et naturel exceptionnel, abritant de nombreux sites classés au patrimoine mondial de l'UNESCO. Cependant, l'accès à l'information et la réservation simplifiée pour ces destinations restent sous-optimaux pour les touristes nationaux et internationaux.

\textbf{WEJHA} (وجهة) est une plateforme moderne de découverte et de réservation de visites touristiques sous forme de MVP (Minimum Viable Product). Elle vise à centraliser les informations sur les sites phares des 58 Wilayas d'Algérie, en proposant :
\begin{itemize}
    \item Un catalogue enrichi de sites touristiques avec filtrage avancé.
    \item Un système de réservation de créneaux avec gestion d'affluence.
    \item Une expérience utilisateur (UX) moderne et réactive adaptée aux mobiles.
    \item Des fonctionnalités innovantes telles qu'un chatbot d'IA de voyage et un mini-jeu de quiz patrimonial.
\end{itemize}

\subsection{Public Cible}
\begin{enumerate}
    \item \textbf{Touristes Locaux :} Résidents souhaitant explorer d'autres wilayas durant les week-ends ou vacances.
    \item \textbf{Diaspora et Touristes Étrangers :} Visiteurs à la recherche d'itinéraires et d'informations fiables.
    \item \textbf{Opérateurs Touristiques :} Guides et gestionnaires de sites pour la régulation des flux de visiteurs.
\end{enumerate}

% --- Section 2: Spécifications Fonctionnelles ---
\section{Spécifications Fonctionnelles}

\subsection{Base de Données des Sites Touristiques}
Le MVP intègre un catalogue initial de plus de 85 sites touristiques répartis sur l'ensemble du territoire national algérien. Chaque site est caractérisé par :
\begin{itemize}
    \item \textbf{Identifiant unique (ID) :} Numérique.
    \item \textbf{Nom :} Ex. \textit{Casbah d'Alger}, \textit{Tassili n'Ajjer}, \textit{Pont de Sidi M'Cid}.
    \item \textbf{Wilaya \& Numéro :} Liés à la nomenclature officielle (Ex. Alger - 16, Constantine - 25, Djanet - 56).
    \item \textbf{Catégories multiples :} Historique, Nature, Sahara, Religieux, UNESCO, Aventure, Culture, Architecture, Plage, Vue Panoramique, Luxe, Famille, Neige, etc.
    \item \textbf{Prix d'entrée (DZD) :} Tarification locale.
    \item \textbf{Évaluation :} Note moyenne sur 5.0.
    \item \textbf{Durée moyenne de visite :} Exprime en heures ou jours.
    \item \textbf{Niveau d'accessibilité :} Accessible, Moyenne, Difficile (pour randonnées ou désert).
\end{itemize}

\subsection{Système de Filtrage Avancé}
Afin de faciliter la recherche parmi le large catalogue, la plateforme intègre un composant de filtrage multicritères fonctionnant en temps réel :
\begin{table}[h]
    \centering
    \caption{Structure du système de filtrage}
    \vspace{0.2cm}
    \begin{tabularx}{\textwidth}{lX}
        \toprule
        \textbf{Type de Filtre} & \textbf{Options / Comportement} \\
        \midrule
        \textbf{Catégories} & Sélection multiple parmi les 15+ catégories du patrimoine. \\
        \textbf{Gamme de Prix} & Sélecteur de prix minimal et maximal (de 0 DZD à 5000 DZD). \\
        \textbf{Évaluation (Rating)} & Filtrage par note minimale (ex. $\ge$ 4.0 étoiles). \\
        \textbf{Accessibilité} & Choix exclusif (Accessible, Moyenne, Difficile). \\
        \textbf{Durée} & Filtrage selon le type d'excursion (Court séjour, Journée, Multi-jours). \\
        \bottomrule
    \end{tabularx}
\end{table}

\subsection{Assistant de Voyage IA (Wejha AI)}
Intégré sous forme de widget flottant interactif, l'assistant \textbf{Wejha AI} exploite l'API Google Gemini pour guider les utilisateurs. Ses capacités incluent :
\begin{itemize}
    \item \textbf{Recommandations contextuelles :} Analyse des profils de voyageurs sélectionnés (ex. \textit{«~Budget < 500 DZD~»}, \textit{«~UNESCO \& Histoire~»}, \textit{«~Aventure Sahara~»}).
    \item \textbf{Chat interactif :} Réponses en temps réel sur les conseils de voyage en Algérie (périodes idéales pour le Grand Sud, transport, formalités).
    \item \textbf{Redirection interactive :} Les noms de sites mentionnés par l'IA sous la forme \textbf{**Nom du Site**} sont cliquables et ouvrent directement la fiche du site sur la plateforme.
\end{itemize}

\subsection{Gamification : Le Défi Patrimoine}
Un jeu-quiz interactif est intégré pour sensibiliser les utilisateurs au patrimoine algérien et augmenter le taux de rétention :
\begin{itemize}
    \item \textbf{Structure :} 5 questions tirées au sort sur l'histoire et la géographie des monuments nationaux.
    \item \textbf{Apprentissage actif :} Chaque question affiche une anecdote descriptive (\textit{«~Le saviez-vous ?~»}) après validation.
    \item \textbf{Grades décernés :}
        \begin{itemize}
            \item \textbf{5/5 :} Ambassadeur du Patrimoine 🇩🇿
            \item \textbf{3-4/5 :} Explorateur Éclairé
            \item \textbf{0-2/5 :} Voyageur Curieux
        \end{itemize}
\end{itemize}

\subsection{Module de Réservation et Gestion d'Affluence}
Pour chaque site, l'utilisateur peut planifier sa visite :
\begin{enumerate}
    \item \textbf{Choix de la date et du créneau horaire} (Matin, Après-midi, Soir).
    \item \textbf{Indicateur d'affluence en temps réel} (\textit{Faible}, \textit{Modérée}, \textit{Forte}, \textit{Complet}) calculé pour réguler les visites de sites fragiles (ex. Casbah ou musées).
    \item \textbf{Moyens de Paiement acceptés :} Cartes locales interbancaires (CIB, Edahabia) et option de paiement en espèces sur site (Cash).
\end{enumerate}

% --- Section 3: Design et Expérience Utilisateur ---
\section{Design et Identité Visuelle}
La plateforme applique une esthétique moderne et épurée combinant des couleurs inspirées des paysages algériens (mer, oasis, dunes et artisanat traditionnel).

\subsection{Charte Graphique (Design Tokens)}
\begin{itemize}
    \item \textbf{Couleur Primaire (Vert Oasis) :} \texttt{\#0F6E56} (utilisé pour les boutons principaux, en-têtes et éléments actifs).
    \item \textbf{Couleur Secondaire (Vert Menthe) :} \texttt{\#1D9E75} (utilisé pour les accents et les états de survol).
    \item \textbf{Couleur d'Accent (Ambre Désert) :} \texttt{\#BA7517} (utilisé pour les notes, prix et distinctions).
    \item \textbf{Arrière-plan (Sable Clair) :} \texttt{\#F8F5EF} (donne une texture chaleureuse et naturelle).
    \item \textbf{Surface (Blanc Pur) :} \texttt{\#FFFFFF} (utilisé pour les cartes et conteneurs).
\end{itemize}

\subsection{Typographie et Espacement}
\begin{itemize}
    \item \textbf{Polices :} \textit{DM Sans} pour le corps de texte et l'interface utilisateur générale ; \textit{Playfair Display} pour les titres majeurs et invitations à l'exploration.
    \item \textbf{Bordures :} Coins arrondis généreux (\texttt{12px} à \texttt{24px}) pour adoucir l'interface et évoquer l'architecture vernaculaire.
\end{itemize}

\subsection{Optimisation Mobile-First}
L'interface a été conçue en priorité pour les téléphones intelligents en raison de l'usage prédominant en situation de mobilité :
\begin{itemize}
    \item \textbf{Navigation Basse (MobileBottomNav) :} Barre fixe en bas de l'écran avec 5 onglets d'accès rapide (Accueil, Catalogue, Favoris, Défi, IA).
    \item \textbf{Menu Latéral coulissant (MobileNav) :} Se déploie de la droite pour l'accès aux paramètres de profil et à la connexion.
\end{itemize}

% --- Section 4: Architecture Technique ---
\section{Architecture Technique et Code}

\subsection{Pile Technologique (Stack)}
\begin{itemize}
    \item \textbf{Frontend :} React 18 / 19 avec TypeScript pour un typage statique robuste.
    \item \textbf{Build Tool :} Vite pour des performances optimales en développement et en production.
    \item \textbf{Animations :} Framer Motion / Motion.react pour des transitions fluides et micro-interactions tactiles.
    \item \textbf{Styles :} CSS Vanilla et Tailwind CSS pour une conception adaptative.
    \item \textbf{Icônes :} Lucide React.
\end{itemize}

\subsection{Structure des Composants Majeurs}
L'application est structurée de manière modulaire autour de composants réutilisables clés :
\begin{lstlisting}[caption={Exemple de typage pour les sites touristiques}]
type SiteCategory = 'Historique' | 'Nature' | 'Sahara' | 'Religieux' | 'UNESCO' | 'Aventure' | 'Culture' | 'Architecture' | 'Plage' | 'Vue Panoramique' | 'Luxe' | 'Famille' | 'Neige' | 'Vie Nocturne' | 'Musée' | 'Randonnée' | 'Monument' | 'Jardin';

interface Site {
  id: number;
  name: string;
  wilaya: string;
  wilayaNumber: number;
  categories: SiteCategory[];
  price: number;
  rating: number;
  duration: string;
  accessibility: string;
  image: string;
  description: string;
}
\end{lstlisting}

\subsection{Composants UI développés}
\begin{itemize}
    \item \textbf{\texttt{AdvancedFilters.tsx} :} Gère l'état et l'interface utilisateur des filtres complexes avec un indicateur du nombre de filtres actifs.
    \item \textbf{\texttt{SiteCard.tsx} :} Affiche les sites en deux variantes : \texttt{grid} (vue compacte de catalogue) et \texttt{featured} (mise en avant grand format dans le feed).
    \item \textbf{\texttt{SkeletonLoader.tsx} :} Réduit le temps de chargement perçu grâce à des squelettes animés avec effet de miroitement (\textit{shimmer effect}).
    \item \textbf{\texttt{ErrorBoundary.tsx} :} Intercepte les erreurs de rendu React au niveau des composants et propose un écran de secours personnalisé avec bouton de réinitialisation.
\end{itemize}

% --- Section 5: Exigences Non-Fonctionnelles ---
\section{Exigences Non-Fonctionnelles}

\subsection{Résilience et Gestion des Erreurs}
\begin{tcolorbox}[colback=red!5,colframe=red!50!black,title=Zéro Écran Blanc (White Screen of Death)]
    Toutes les routes critiques et les composants consommant des données asynchrones (comme l'assistant IA et le flux d'actualités) doivent être enveloppés dans un \texttt{ErrorBoundary}. L'application ne doit jamais crasher globalement à cause d'une panne d'API.
\end{tcolorbox}

\subsection{Accessibilité (WCAG 2.1 AA)}
\begin{itemize}
    \item \textbf{Contraste :} Respect des ratios de contraste de couleur pour les malvoyants (minimum 4.5:1).
    \item \textbf{Navigation au Clavier :} Prise en charge des focus visibles pour tous les boutons personnalisés (\texttt{Button.tsx}) et éléments interactifs.
    \item \textbf{Lecteurs d'écran :} Balises sémantiques HTML5 (\texttt{<header>}, \texttt{<nav>}, \texttt{<main>}, \texttt{<footer>}) et attributs \texttt{aria-label} appropriés.
\end{itemize}

\subsection{Performance}
\begin{itemize}
    \item \textbf{Temps de Chargement Perçu :} Utilisation systématique de squelettes de chargement (\textit{Skeletons}) lors des appels d'API afin d'abaisser la perception d'attente de 2 à 3 secondes.
    \item \textbf{Optimisation des Images :} Redimensionnement dynamique via des URL distantes optimisées (\texttt{unsplash/pexels}) utilisant les paramètres de largeur \texttt{w=800}.
\end{itemize}

\end{document}
